\Article{4}
Apart from the magnetic forces we see individual bodies exert on one another when they are close enough together, another force acts on the magnetic fluids. Since it manifests itself everywhere on Earth, we attribute it to the globe itself and call it terrestrial magnetism. This force shows itself in two ways. If bodies of the second kind, in which the magnetism has developed, are supported at their center of gravity, the force tends to turn them in a definite direction. In bodies of the first kind, by contrast, this force spontaneously separates the magnetic fluids. The separation can be made quite perceptible by choosing bodies of a suitable shape and placing them in a suitable position. We explain both phenomena by supposing that at any given place this force drives the northern magnetic fluid in a definite direction and the southern fluid with equal intensity in the opposite direction. It is the former direction that we always mean by the direction of terrestrial magnetism. It is therefore determined by its inclination to the \SourceLeaf{93}horizontal plane and by the declination of the vertical plane in which it acts from the meridian plane. That vertical plane is called the \textit{magnetic meridian plane}. The intensity of terrestrial magnetism is to be measured by the motive force it exerts on a unit of free magnetic fluid.

This force differs from place to place on Earth and also varies at the same place over centuries and years. It varies as well through the summers of the year and the hours of the day. The variation in direction has long been known. So far, variation in intensity could be observed only over the hours of the day, since we have lacked the means to observe it over longer intervals. Expressing the intensity in absolute units will remedy this deficiency in the future.

\Article{5}
To calculate the action of terrestrial magnetism on magnetic bodies of the second kind, which we shall always mean from now on, imagine such a body divided into infinitely small parts. Let $\mathrm{d}m$ be the element of free magnetism in a particle whose coordinates with respect to three mutually perpendicular planes fixed relative to the body are denoted by $x,y,z$. We take elements of the southern fluid to be negative. First, then, it is clear that the integral $\int\mathrm{d}m$ taken over the entire body, or indeed over any measurable part of it, is $=0$. Set $\int x\mathrm{d}m=X$, $\int y\mathrm{d}m=Y$, $\int z\mathrm{d}m=Z$. We may call these quantities the moments of free magnetism with respect to the three fundamental planes, or with respect to the axes perpendicular to them. If $a$ denotes an arbitrary constant, then $\int(x-a)\mathrm{d}m=X$, so the moment with respect to a given axis clearly depends only on its direction and is independent of its origin. If we draw a fourth axis through the origin of the coordinates, making angles $A,B,C$ with the original axes, the moment of the element $\mathrm{d}m$ with respect to this axis will be $=(x\cos A+y\cos B+z\cos C)\mathrm{d}m$, and hence the moment of free magnetism in the whole body will be
\[
=X\cos A+Y\cos B+Z\cos C=V
\]
Set
\[
\BareRoot{(XX+YY+ZZ)}=M,\enspace\text{and }X=M\cos\alpha,\enspace Y=M\cos\beta,\quad Z=M\cos\gamma
\]
and draw a fifth axis making angles $\alpha,\beta,\gamma$ with the original three and an angle $\omega$ with the fourth. Since we know that
\SourceLeaf{94}
\[
\cos\omega=\cos A\cos\alpha+\cos B\cos\beta+\cos C\cos\gamma,\quad\text{we get }V=M\cos\omega
\]
We call this fifth axis simply the \textit{magnetic axis} of the body and take its \textit{direction} to correspond to the positive value of the root $\BareRoot{(XX+YY+ZZ)}$. If the fourth axis coincides with this magnetic axis, the moment $V$ becomes $=M$, which is clearly the greatest of all the moments. The moment with respect to any other axis is found by multiplying this greatest moment, which we may simply call the moment of magnetism whenever there is no risk of ambiguity, by the cosine of the angle between that axis and the magnetic axis. The moment with respect to any axis perpendicular to the magnetic axis is $=0$, and the moment with respect to any axis that makes an obtuse angle with the magnetic axis is negative.

The magnetic axis is thus no single definite straight line. It can be drawn through any point and specifies only a direction. In other words, there are infinitely many magnetic axes parallel to one another. If we choose any one of these at will and give it a definite length, its ends are called poles. The axis is directed from the southern pole toward the northern one.

\Article{6}
If a force of constant intensity and direction acts on the individual particles of the magnetic fluids, the resulting total force on the body is easily derived from the principles of statics. In the bodies we are considering, these particles have, as it were, lost their fluidity and form a single rigid mass with the weighable body. Let a motive force act on any magnetic molecule $\mathrm{d}m$, with value $=P\mathrm{d}m$ in direction $D$. For molecules of the southern fluid, the negative sign itself already implies the opposite direction. Let $A,B$ be two points of the body lying in the direction of the magnetic axis, and let their distance be $=r$, taken as positive when the magnetic axis runs from $A$ toward $B$. We readily see that if we add two new forces, each $=\frac{PM}{r}$, one acting at $A$ in direction $D$ and the other at $B$ in the opposite direction, all these forces will be in equilibrium. The original forces are therefore equivalent to two forces $=\frac{PM}{r}$, one acting at $B$ in direction $D$ and the other at $A$ in the opposite direction. Clearly these two forces cannot be combined into a single force.

If, in addition to the force $P$, another similar force $P'$ acts on the magnetic fluids of the body in direction $D'$, we may again substitute two others for it. They may act at the same points $A,B$ or, more \SourceLeaf{95}generally, at other points $A',B'$, provided $A'B'$ is also a magnetic axis. If we set the distance $A'B'=r'$, these forces must be $=\frac{P'M}{r'}$, acting at $B'$ in direction $D'$ and at $A'$ in the opposite direction. The same applies to several forces.

Within a space as small as the body to be experimented on occupies, we can safely take the intensity and direction of the terrestrial magnetic force to be constant throughout, even though they may vary with time. We may therefore apply to it what we have just said. It can be convenient, however, to resolve it at once into two forces, a horizontal one $=T$ and a vertical one $=T'$, which points downward in our regions. If we wish to replace the latter by two other forces acting at points $A',B'$, we can choose both the point $A'$ and the distance $A'B'=r'$ at will. We shall take the center of gravity for $A'$ and, denoting the body's weight, that is, the motive force gravity exerts on its mass, by $p$, set $\frac{T'M}{p}=r'$. In this way the effect of the force $T'$ is resolved into an upward force $=p$ at $A'$ and an equal downward force at $B'$. Gravity itself clearly cancels the former, so the effect of the vertical terrestrial magnetic force comes down to a shift of the center of gravity from $A'$ to $B'$. Clearly, in regions where the terrestrial magnetic force makes an acute angle with the vertical line, or where its vertical part drives the northern magnetic fluid upward, a similar shift of the center of gravity occurs along the magnetic axis toward the southern pole.

This way of looking at the matter makes it evident that whatever experiments we make with a magnetic needle in a \textit{single} magnetic state, we cannot infer the inclination from them alone. We must already know the position of the \textit{true} center of gravity by other means. This position is usually established before magnetizing the needle. But the method is hardly reliable, since a steel needle generally acquires some magnetism, however weak, during manufacture. To determine the inclination, then, we must induce \textit{another} shift of the center of gravity by suitably altering the magnetic state of the needle. To make this shift as different as possible from the first, we must reverse the poles, which can give us twice the displacement. Even for needles of the most suitable dimensions and magnetized to saturation, the shift of the center of gravity cannot exceed a certain limit. For a single shift this limit is about $0{,}4$ millimeters in our regions, and it \SourceLeaf{96}remains below $0{,}6$ millimeters in regions where the vertical force is greatest. This also shows how much mechanical precision is required in needles intended to determine the inclination.

\Article{7}
Suppose that some point $C$ of a magnetic body is fixed. For equilibrium it is necessary and sufficient that the plane drawn through $C$, the center of gravity and the magnetic axis coincide with the magnetic meridian plane, and that the moments with which the terrestrial magnetic force and gravity tend to turn this plane about $C$ cancel one another. The latter condition amounts to the following. Let $T$ denote the horizontal part of the terrestrial magnetic force and $i$ the inclination of the magnetic axis to the horizontal plane. Then $TM\sin i$ must equal the product of the body's weight and the distance of the displaced center of gravity $B'$ from the vertical line through $C$. Clearly this distance must lie on the southern or northern side according to whether $i$ is an elevation or a depression, and for $i=0$, $B'$ lies on that vertical line itself. Now suppose the body has been turned about this vertical line until the magnetic axis lies in a vertical plane whose magnetic azimuth, that is, its angle with the northern part of the magnetic meridian, is $=u$. We can take either eastward or westward angles as positive at will. Terrestrial magnetism will exert a force tending to turn the body about the vertical axis, that is, to decrease the angle $u$. Its moment will be $=TM\cos i\sin u$, and the body will oscillate about this axis. The duration of the oscillations can be calculated by known methods. Specifically, let $K$ denote the body's moment of inertia about the axis of oscillation, that is, the sum of the weighable molecules multiplied by the squares of their distances from the axis, and let $\pi$ denote, as usual, the semicircumference of a circle of radius $=1$. The time of one infinitely small oscillation will be $=\pi\BareRoot{\frac{K}{TM\cos i}}$. This holds if the unit of accelerating force underlying the quantities $T,M$ is the one that produces a velocity $=1$ in a unit of time. The reduction from finite to infinitely small oscillations can be calculated in much the same way as for the oscillations of a pendulum. If, then, the time of one infinitely small oscillation has been found from observations to be $=t$, we have $TM=\frac{\pi\pi K}{tt\cos i}$. Hence, if the body is suspended with its magnetic axis horizontal, as we shall always assume from now on, we have
\[
TM=\frac{\pi\pi K}{tt}
\]
If we prefer to take gravity as the unit of accelerating force, we would have to divide this \SourceLeaf{97}value by $\pi\pi l$, where $l$ denotes the length of a simple pendulum that makes one swing per unit of time. Thus in general we have $TM=\frac{K}{ttl\cos i}$, or in our case $TM=\frac{K}{ttl}$.
